\documentclass[12pt]{article}
\usepackage[margin=1in]{geometry}
\usepackage{amsmath}
\usepackage{hyperref}
\hypersetup{hidelinks}
\begin{document}

\begin{center}
{\Large\textbf{Nature Research Reporting Summary}}\\[0.5cm]
{\large Manuscript: ``Sample-Size Phase Transition Ends the Hard-versus-Soft Clustering Debate''}\\[0.3cm]
Corresponding Author: Shuaidong Gao (gsd3247186514@gmail.com)
\end{center}

\vspace{0.5cm}

\section*{Statistics}

\begin{itemize}
\item \textbf{Sample sizes:} TCGA cross-cancer comparison: 33 cancer types ($n = 45$--$1218$ samples; 10 seeds at $d=200$, Extended Data Table~1). Within-cancer resampling: 8 cancers, each resampled at $n \in \{50, 100, 200, 400\}$ (5 seeds). Cross-modal replication: MNIST $n \in \{50$--$3000\}$ (12 sizes, 3 seeds each); PBMC 10k $n \in \{50$--$10000\}$ (15 sizes, 3 seeds each); 20 Newsgroups $n \in \{50$--$10000\}$ (15 sizes, 3 seeds each). Methylation: 8 cancer types, $n = 99$--$554$ (3 seeds). No statistical method was used to predetermine sample size; all available data were used.

\item \textbf{Statistical tests:} Spearman rank correlation (two-sided) was the primary test for monotonic relationships between sample size and method advantage. Mann-Whitney $U$ tests compared small-sample ($n < 323$) vs.\ large-sample ($n \geq 323$) cancer groups. Wilcoxon signed-rank tests evaluated paired CAGate-vs-SSCAGate comparisons. Binomial tests assessed all-$\Delta$-positive claims across the 33-cancer panel. Curve fitting used Levenberg-Marquardt optimization.

\item \textbf{Data exclusions:} No cancer types were excluded from the cross-cancer comparison. The number of top variable genes was capped at $d = \min(200, n-10)$ per cancer to ensure identifiability ($d < n$); for all cancers with $n \geq 210$ this yields $d=200$. For methylation gene aggregation, genes with $<3$ CpG probes were excluded (standard practice).

\item \textbf{Replication:} The phase transition was replicated in: (i) 33 independent TCGA cancer types, (ii) 8 within-cancer resampling experiments, (iii) PBMC single-cell RNA-seq datasets (3k, 6k, 10k), (iv) MNIST image data, (v) 20 Newsgroups text data, (vi) 8 methylation datasets. All replication attempts confirmed the phase transition.

\item \textbf{Randomization:} For all subsampling experiments, random seeds (42 for MNIST, 42 for PBMC 10k, 43 for 20 Newsgroups) were used to select cell/image/document subsets. For CAGate/SSCAGate runs: 10 seeds per cancer for the $d=200$ cross-cancer comparison (Extended Data Table~1), 5 seeds for within-cancer resampling (Figure~1b), and 3 seeds for cross-modal replication. K-means initialization was randomized.

\item \textbf{Blinding:} Not applicable. All analyses were reproducible given fixed random seeds.
\end{itemize}

\section*{Experimental Design}

\begin{itemize}
\item \textbf{Study design:} This is a computational study with three independent experimental paradigms: (i) cross-cancer comparison of 33 independent datasets, (ii) within-cancer controlled resampling, and (iii) cross-modal replication on completely unrelated data types. No treatment groups or interventions were involved.

\item \textbf{Controls:} Standard NOTEARS (without gating) served as the baseline in all experiments. Synthetic data (Erd\H{o}s-R\'{e}nyi DAGs) were used as a negative control. Random cluster assignment was used as an ablation control.

\item \textbf{Data sources:} All data are publicly available. TCGA gene expression: GDC Data Portal and UCSC Xena. Single-cell RNA-seq: 10x Genomics and Scanpy. MNIST: Yann LeCun's website. 20 Newsgroups: public NLP corpus. DNA methylation: Firehose Broad GDAC.
\end{itemize}

\section*{Software and Algorithms}

\begin{itemize}
\item \textbf{Programming languages:} Python 3.12
\item \textbf{Key packages:} PyTorch 2.11, scikit-learn 1.6, Scanpy 1.10, NumPy 1.26, SciPy 1.11
\item \textbf{Custom code:} CDSM framework (v3.0.0), available at \url{https://github.com/gsd3247186514-cell/SSCAGate-Nature} (MIT license)
\item \textbf{Algorithms:} CAGate (two-stage cluster-aware gradient gating), SSCAGate (self-supervised causal gating with end-to-end cluster optimization), NOTEARS (Zheng et al., 2018)
\item \textbf{Hardware:} AMD Ryzen 9 8945HX (16-core CPU), NVIDIA RTX 5060 Laptop GPU (8GB), 16GB DDR5 RAM, Windows 11
\end{itemize}

\section*{Data and Code Availability}

\begin{itemize}
\item \textbf{Data availability statement:} Included in the main manuscript. All datasets are publicly accessible. Full replication package at \url{https://github.com/gsd3247186514-cell/SSCAGate-Nature}.
\item \textbf{Code availability statement:} Included in the main manuscript. Full source code at \url{https://github.com/gsd3247186514-cell/SSCAGate-Nature}.
\end{itemize}

\section*{Competing Interests}

A Chinese patent application covering the Cluster-Aware Gating mechanism has been filed (May 2026). The author declares no other competing financial or non-financial interests.

\end{document}
